\section{Exact-value identification}\label{sec:identification}
\begin{theorem}[Identification of the canonical supremum]\label{thm:O7}
The finite-dimensional projective tensor supremum is the unique matching
parameter: \(\St=s_*\).
\end{theorem}
\begin{proof}
The inference uses precisely the same interval, normalization, recurrence,
ratio coordinate and intrinsic matching objects throughout:
\[
\begin{array}{rcll}
 \St\in I
 &&&\text{Proposition~\ref{prop:R1}},\\
 \mathcal C_{\St}\ne\varnothing
 &&&\text{Theorem~\ref{thm:R0}},\\
 \mathcal M_{\St}\ne\varnothing
 &&&\text{Proposition~\ref{prop:bridge} at }s=\St,\\
 \St\in\pi_s(\mathcal M)
 &&&\text{definition of the projection},\\
 \pi_s(\mathcal M)=\{s_*\}
 &&&\text{Theorem~\ref{thm:O6}}.
\end{array}
\]
Membership in the singleton gives \(\St=s_*\). Proposition
\ref{prop:scalar} gives the equivalent residual characterization.
\end{proof}
This composition uses one carrier at the true value and uniqueness of the
parameter. It does not require uniqueness of a carrier or optimizer, a
finite-dimensional attaining strategy, any assertion labelled R2, or
commuting=tensor equality.
